\chapter{Reference Frames and Frame Conversions}
\label{chap:frame-conversions}

This chapter specifies every reference frame \lupnt{} knows about, the exact
rotation and translation that relates each adjacent pair, the time scale each
rotation is evaluated in, and the velocity transformation that accompanies
each rotation. The public dispatch --- \cpp{ConvertFrame},
\cpp{GetFrameRotationTranslation}, \cpp{GetFrameRotationTranslationRv} and
\cpp{GetFrameCenter} --- is implemented in
\file{cpp/lupnt/conversions/frame_converter.cc} and declared in
\file{cpp/lupnt/conversions/frame_converter.h}, which also owns the
\cpp{Frame} enumeration. The individual body-pair rotations
(\cpp{GcrfToItrf}, \cpp{GcrfToMoonCi}, \cpp{RotMoonCiToPa}, the IAU planet
orientation, \dots) live in
\file{cpp/lupnt/conversions/frame_conversions.cc} and
\file{cpp/lupnt/conversions/frame_conversions.h}. The topocentric
(east--north--up, azimuth--elevation--range) helpers are in
\file{cpp/lupnt/conversions/coordinate_conversions.cc}.

The Earth models follow the IAU 2006/2000A conventions of
\citep{iers2010}, with the classical background of \citep{montenbruck2000}
and \citep{vallado2013}. The Earth--Moon translation and the lunar mantle
libration angles come from the JPL DE440 ephemeris \citep{park2021de440}.
The optional high-accuracy path is calibrated against
SPICE \citep{acton1996spice}. The cislunar frame conventions and their
navigation consequences are developed in \citep[Ch.~2]{iiyama2026thesis}.
Time-scale conversions themselves --- $\TDB\to\TT$, $\TDB\to\UTC$,
$\TDB\to\UTone$ --- are specified in \cref{chap:time-conversions} and are
used here as black boxes.

\section{Epoch, state and unit contract}
\label{sec:fc-contract}

Every frame-conversion entry point takes its epoch as \emph{barycentric
dynamical time in seconds past the J2000 epoch}:
%
\begin{equation}
  \label{eq:fc-epoch}
  t \equiv t_{\TDB}
  \quad [\si{\second}\ \text{past J2000}].
\end{equation}
%
A caller holding an epoch in any other scale must convert it first
(\cref{chap:time-conversions}); no frame function inspects or repairs the
scale of its argument. Internally the Earth-orientation routines re-derive
whichever scale each physical effect depends on:
$\TT$ for the precession--nutation series and the TIO locator, $\UTone$ for
the Earth Rotation Angle, and $\UTC$ for the tabulated Earth orientation
parameters (EOP). \Cref{tab:fc-timescales} collects the dependencies.

\begin{table}[H]
  \centering
  \caption{Time scale consumed by each elementary rotation. All are derived
    internally from the single $t_{\TDB}$ argument.}
  \label{tab:fc-timescales}
  \begin{tabularx}{\textwidth}{l l l L}
    \toprule
    Rotation & Implementing function & Scale & Reason \\
    \midrule
    $\mat{R}_{\mathrm{PN}}$    & \cpp{RotPrecessionNutation} & $\TT$   & the IAU 2006/2000A $(X,Y,s)$ series is tabulated against Julian date TT \\
    $\mat{R}_{\mathrm{ERA}}$   & \cpp{RotSideralMotion}      & $\UTone$ & the Earth Rotation Angle is defined as a linear function of UT1 \\
    $\dot{\mat{R}}_{\mathrm{ERA}}$ & \cpp{RotSideralMotionDot} & $\UTone$, $\UTC$ & rate from EOP length-of-day (sampled at UTC) or from the fitted $\UTone-\UTC$ \\
    $\mat{R}_{\mathrm{polar}}$ & \cpp{RotPolarMotion}        & $\TT$, $\UTC$ & TIO locator $s'$ in TT; $(x_p,y_p)$ interpolated from the EOP table at UTC \\
    $\mat{B}_e$                & \cpp{RotGcrfToEme}          & ---     & constant frame bias, no epoch argument \\
    $\mat{R}_{\mathrm{PA},\mathrm{CI}}$ & \cpp{RotMoonCiToPa} & $\TDB$  & DE440 libration Chebyshev coefficients are indexed by TDB \\
    $\mat{B}_{\mathrm{moon}}$  & \cpp{RotMoonPaToMe}         & ---     & constant PA-to-ME bias, no epoch argument \\
    $\mat{R}_{\mathrm{OP}}$    & \cpp{RotMoonOpToCi}         & $\TDB$  & built from the DE440 Moon--Earth state \\
    $\mat{R}_{\mathrm{FIXED},\mathrm{CI}}$ & \cpp{RotBodyCiToFixed} & $\TT$ & IAU orientation elements are linear in days/centuries past J2000 TT \\
    \bottomrule
  \end{tabularx}
\end{table}

\paragraph{Units.} Frame conversions are unit-agnostic \emph{for the state}:
they apply a dimensionless rotation and add a translation obtained from the
ephemeris, so the output carries whatever length unit the input carried
\emph{provided} the ephemeris is queried in the same unit. \lupnt{}'s
ephemeris layer returns \si{\meter} and \si{\meter\per\second}, so the
documented contract of every \cpp{ConvertFrame} overload is
\si{\meter} and \si{\meter\per\second}. Angles are radians throughout,
except in the topocentric helpers of \cref{sec:fc-topocentric}, where
latitude and longitude are in degrees. The epoch is always seconds.

\begin{lupntwarning}
  Mixing units silently produces a wrong answer for any conversion that
  \emph{translates} (ICRF$\leftrightarrow$GCRF,
  GCRF$\leftrightarrow$MOON\_CI, the planet frames): a state in kilometres
  will have a metre-valued ephemeris offset added to it. Pure rotations
  (GCRF$\leftrightarrow$ITRF, GCRF$\leftrightarrow$EME,
  MOON\_CI$\leftrightarrow$MOON\_PA, MOON\_PA$\leftrightarrow$MOON\_ME) are
  unaffected.
\end{lupntwarning}

\section{The frame set}
\label{sec:fc-frames}

\Cref{tab:fc-frames} enumerates every value of the \cpp{Frame} enum declared
in \file{cpp/lupnt/conversions/frame_converter.h}, its origin, the definition
of its axes, and whether \cpp{ConvertFrameBase} implements a route to it.
Two aliases exist at the enum level: \code{ECEF} is a synonym for
\code{ITRF} and \code{ECI} is a synonym for \code{EME}, so they compare
equal and cannot be distinguished at runtime.

\begin{table}[H]
  \centering
  \caption{The \cpp{Frame} enumeration: origin, axis definition, and
    implementation status.}
  \label{tab:fc-frames}
  \small
  \begin{tabularx}{\textwidth}{l l L l}
    \toprule
    \code{Frame} & Origin & Axes & Status \\
    \midrule
    \code{UNDEFINED} & --- & sentinel for an untagged state & throws \\
    \code{ICRF}   & solar-system barycenter & ICRS axes: the kinematically non-rotating extragalactic realization & hub \\
    \code{GCRF}   & Earth center of mass & ICRS axes, geocentric origin & hub \\
    \code{ITRF} (\code{ECEF}) & Earth center of mass & IERS terrestrial pole and prime meridian; co-rotates with the crust & implemented \\
    \code{EME} (\code{ECI}) & Earth center of mass & mean equator and equinox of J2000 & implemented \\
    \code{EMR}    & Earth--Moon barycenter & synodic: $\uvec{x}$ along Earth$\to$EMB, $\uvec{z}$ along the EMB orbit normal, $\uvec{y}$ completing the in-plane triad & implemented \\
    \code{SER}    & Earth & Sun--Earth rotating & declared only \\
    \code{GSE}    & Earth & geocentric solar ecliptic & declared only \\
    \code{MOD}    & Earth & mean equator and equinox of date & declared only \\
    \code{TOD}    & Earth & true equator and equinox of date & declared only \\
    \code{MOON\_CI} & Moon center of mass & ICRS/GCRF axes, selenocentric origin & hub \\
    \code{MOON\_PA} & Moon center of mass & DE440 lunar principal axes (mantle frame) & implemented \\
    \code{MOON\_ME} & Moon center of mass & mean-Earth/polar axes; a fixed bias from PA & implemented \\
    \code{MOON\_OP} & Moon center of mass & orbit plane: $\uvec{z}$ along the Moon--Earth orbit normal, $\uvec{x}$ the ME pole projected into that plane & implemented \\
    \code{<PLANET>\_CI} & planet center & ICRS axes, planet-centered; \code{MERCURY} \dots \code{NEPTUNE} & implemented \\
    \code{<PLANET>\_FIXED} & planet center & IAU body-fixed axes (pole $\alpha,\delta$ and prime meridian $W$) & implemented \\
    \bottomrule
  \end{tabularx}
\end{table}

Requesting a conversion to or from \code{SER}, \code{GSE}, \code{MOD},
\code{TOD} or \code{UNDEFINED} makes \cpp{ConvertFrameBase} raise a
\code{std::runtime\_error} carrying the message \code{"Not implemented"};
these values exist so that external data tagged with them round-trips
through the enum, not because a transformation is available.

\paragraph{Frame centers.} The map \cpp{frame\_centers} and its accessor
\cpp{GetFrameCenter} give the natural central body of a frame, used by the
dynamics and clock models to pick a gravitational parameter and a
relativity reference:
%
\begin{align}
  \label{eq:fc-centers}
  \operatorname{center}(\mathrm{ICRF}) &= \mathrm{SSB}, \nonumber\\
  \operatorname{center}(\mathrm{GCRF})
    = \operatorname{center}(\mathrm{ITRF})
    &= \operatorname{center}(\mathrm{EME}) = \mathrm{Earth}, \\
  \operatorname{center}(\mathrm{MOON\_CI})
    = \operatorname{center}(\mathrm{MOON\_PA})
    = \operatorname{center}(\mathrm{MOON\_ME})
    &= \operatorname{center}(\mathrm{MOON\_OP}) = \mathrm{Moon},
    \nonumber
\end{align}
%
together with $\operatorname{center}(\code{<PLANET>\_CI}) =
\operatorname{center}(\code{<PLANET>\_FIXED}) = \text{that planet}$.

\begin{lupntnote}
  \code{EMR} has no entry in \cpp{frame\_centers}, so
  \cpp{GetFrameCenter}\code{(Frame::EMR)} throws even though
  \cpp{ConvertFrame} handles \code{EMR} states. Code that needs a center for
  a rotating frame must special-case it.
\end{lupntnote}

\section{The affine transform}
\label{sec:fc-affine}

\subsection{Positions}

For position-only conversions, the mapping from frame $A$ to frame $B$ is
affine:
%
\begin{equation}
  \label{eq:fc-affine}
  \vec{r}_B = \mat{R}_{BA}(t)\,\vec{r}_A + \vec{p}_{BA}(t),
\end{equation}
%
where $\mat{R}_{BA} \in SO(3)$ rotates vector \emph{components} from the $A$
axes into the $B$ axes, and $\vec{p}_{BA}$ is the position of the origin of
$A$ expressed in $B$ axes. \cpp{GetFrameRotationTranslation} returns exactly
the pair $(\mat{R}_{BA}, \vec{p}_{BA})$, and it recovers them without any
knowledge of the routing by transforming the identity basis and the zero
vector:

\begin{lstlisting}[language=cpp, caption={Recovering the affine pair from
  the generic dispatcher (\texttt{cpp/lupnt/conversions/frame\_converter.cc}).},
  label={lst:fc-rt}]
// cpp/lupnt/conversions/frame_converter.cc :: GetFrameRotationTranslation
std::pair<Mat3, Vec3> GetFrameRotationTranslation(Real t_tdb, Frame from_frame, Frame to_frame) {
  MatX3 I3 = Mat3::Identity();
  Mat3 M = ConvertFrame(t_tdb, I3, from_frame, to_frame);          // images of e1,e2,e3
  MatX3 r0 = MatX3::Zero(1, 3);
  Vec3 r = ConvertFrame(t_tdb, r0, from_frame, to_frame).transpose();   // image of 0 = p_BA
  Mat3 R = (M.rowwise() - r.transpose()).transpose();              // subtract the offset
  return std::make_pair(R, r);
}
\end{lstlisting}

\subsection{Cartesian states and the transport term}
\label{sec:fc-transport}

For a Cartesian state $\vec{x} = (\vec{r}, \vec{v})$ the transform is again
affine, now in six dimensions,
%
\begin{equation}
  \label{eq:fc-affine6}
  \begin{bmatrix}\vec{r}_B \\ \vec{v}_B\end{bmatrix}
  =
  \mat{R}_{6,BA}(t)
  \begin{bmatrix}\vec{r}_A \\ \vec{v}_A\end{bmatrix}
  + \vec{p}_{6,BA}(t),
\end{equation}
%
and for a transform composed of a time-dependent rotation and a
time-dependent translation the blocks follow from differentiating
\cref{eq:fc-affine}:
%
\begin{equation}
  \label{eq:fc-vel}
  \vec{v}_B
  = \mat{R}_{BA}\,\vec{v}_A
  + \dot{\mat{R}}_{BA}\,\vec{r}_A
  + \dot{\vec{p}}_{BA},
  \qquad
  \mat{R}_{6,BA}
  = \begin{bmatrix}
      \mat{R}_{BA} & \mat{0} \\
      \dot{\mat{R}}_{BA} & \mat{R}_{BA}
    \end{bmatrix}.
\end{equation}
%
The middle term of \cref{eq:fc-vel} is the transport (or \emph{Coriolis
carry}) contribution. It is worth writing it in its physical form. Because
$\mat{R}_{BA}$ is orthogonal, differentiating
$\mat{R}_{BA}\mat{R}_{BA}\transpose = \mat{I}$ shows that
$\dot{\mat{R}}_{BA}\mat{R}_{BA}\transpose$ is skew-symmetric, so there is a
unique vector $\vec{\omega}_{BA}$ --- the angular velocity of frame $B$ with
respect to frame $A$, expressed in $B$ axes --- with
%
\begin{equation}
  \label{eq:fc-omega}
  \dot{\mat{R}}_{BA}\,\mat{R}_{BA}\transpose = -\left[\vec{\omega}_{BA}\right]_\times,
  \qquad
  \left[\vec{a}\right]_\times \vec{b} \equiv \vec{a}\times\vec{b}.
\end{equation}
%
Substituting into \cref{eq:fc-vel} with $\vec{p}_{BA}=\vec{0}$ gives the
familiar form
%
\begin{equation}
  \label{eq:fc-transport}
  \vec{v}_B
  = \mat{R}_{BA}\,\vec{v}_A - \vec{\omega}_{BA}\times\vec{r}_B .
\end{equation}
%
\lupnt{} never forms $\vec{\omega}_{BA}$ explicitly: it propagates
$\dot{\mat{R}}$ analytically through the product rule, which is both exact
and differentiable by the automatic-differentiation scalar type. The sign in
\cref{eq:fc-omega} follows from \lupnt{}'s \emph{passive} elementary
rotations,
%
\begin{equation}
  \label{eq:fc-elementary}
  \mat{R}_x(\alpha) =
  \begin{bmatrix} 1&0&0\\ 0&\cos\alpha&\sin\alpha\\ 0&-\sin\alpha&\cos\alpha\end{bmatrix},
  \quad
  \mat{R}_y(\alpha) =
  \begin{bmatrix} \cos\alpha&0&-\sin\alpha\\ 0&1&0\\ \sin\alpha&0&\cos\alpha\end{bmatrix},
  \quad
  \mat{R}_z(\alpha) =
  \begin{bmatrix} \cos\alpha&\sin\alpha&0\\ -\sin\alpha&\cos\alpha&0\\ 0&0&1\end{bmatrix},
\end{equation}
%
defined in \file{cpp/lupnt/numerics/math_utils.cc} together with their exact
derivatives $\mat{R}_x'(\alpha,\dot\alpha) = \dd\mat{R}_x/\dd t$ etc.\ (the
\cpp{RotXdot}/\cpp{RotYdot}/\cpp{RotZdot} family). Every composite rotation
in this chapter is built from these three, which is what guarantees that the
reconstructed matrices are orthonormal to floating-point precision no matter
how the underlying angles were obtained.

\cpp{GetFrameRotationTranslationRv} returns
$(\mat{R}_{6,BA}, \vec{p}_{6,BA})$ by the same trick as \cref{lst:fc-rt},
applied to the six Cartesian basis states:

\begin{lstlisting}[language=cpp, caption={Six-state affine transform
  (\texttt{cpp/lupnt/conversions/frame\_converter.cc}).},
  label={lst:fc-rt6}]
// cpp/lupnt/conversions/frame_converter.cc :: GetFrameRotationTranslationRv
MatX6 zero_row = MatX6::Zero(1, 6);
Vec6 t6 = ConvertFrame(t_tdb, zero_row, from_frame, to_frame, false).row(0).transpose();
MatX6 I6 = Mat6::Identity();
MatX6 Y = ConvertFrame(t_tdb, I6, from_frame, to_frame, false);  // 6x6
Mat6 R6;
for (int j = 0; j < 6; ++j) R6.col(j) = Y.row(j).transpose() - t6;
return {R6, t6};
\end{lstlisting}

\subsection{Rotate-only mode}

Every \cpp{ConvertFrame} overload accepts a trailing \code{rotate\_only}
flag. When it is set, the six-state path is replaced by the \emph{three}
dimensional affine pair applied to the position and the bare rotation
applied to the velocity:
%
\begin{equation}
  \label{eq:fc-rotate-only}
  \vec{r}_B = \mat{R}_{BA}\vec{r}_A + \vec{p}_{BA},
  \qquad
  \vec{v}_B = \mat{R}_{BA}\vec{v}_A .
\end{equation}

\begin{lstlisting}[language=cpp, caption={The \texttt{rotate\_only} branch
  (\texttt{cpp/lupnt/conversions/frame\_converter.cc}).},
  label={lst:fc-rotate-only}]
// cpp/lupnt/conversions/frame_converter.cc :: ConvertFrame (Vec6 overload)
Vec6 ConvertFrame(Real t_tdb, const Vec6& rv_in, Frame frame_in, Frame frame_out,
                  bool rotate_only) {
  if (frame_in == frame_out) return rv_in;
  if (!rotate_only) return ConvertFrameBase(t_tdb, rv_in, frame_in, frame_out);
  auto [R, r] = GetFrameRotationTranslation(t_tdb, frame_in, frame_out);
  Vec6 rv_out;
  rv_out.head(3) = R * rv_in.head(3) + r;   // position: rotate AND translate
  rv_out.tail(3) = R * rv_in.tail(3);       // velocity: rotate only
  return rv_out;
}
\end{lstlisting}

\begin{lupntnote}
  The name is slightly misleading: \code{rotate\_only} still translates the
  \emph{position} by $\vec{p}_{BA}$; what it drops is the
  $\dot{\mat{R}}\vec{r} + \dot{\vec{p}}$ contribution to the
  \emph{velocity}. Use it for vector-like quantities --- accelerations,
  line-of-sight directions, local offsets --- where the origin motion must
  not appear. For an absolute position--velocity state, leave it \code{false}.
\end{lupntnote}

\section{The frame graph and the routing algorithm}
\label{sec:fc-graph}

\cpp{ConvertFrameBase} is a recursive \code{switch} over \code{frame\_in}
that hops through three hubs --- ICRF for the solar system, GCRF for the
Earth family, and MOON\_CI for the lunar family --- until \code{frame\_out}
is reached. \Cref{fig:fc-graph} draws the implemented graph.

\begin{figure}[H]
  \centering
  \begin{tikzpicture}[node distance=8mm]
    % --- hubs ---------------------------------------------------------
    \node[tsnode] (icrf)  at (0,0)     {ICRF};
    \node[tsnode] (gcrf)  at (3.4,0)   {GCRF};
    \node[tsnode] (mci)   at (7.0,0)   {MOON\_CI};
    % --- Earth family -------------------------------------------------
    \node[tsnode] (itrf)  at (2.1,1.7)  {ITRF};
    \node[tsnode] (eme)   at (4.7,1.7)  {EME};
    \node[tsnode] (emr)   at (3.4,-1.7) {EMR};
    \node[tsnode] (nyi)   at (0.7,-1.7) {\footnotesize MOD/TOD/SER/GSE};
    % --- Moon family --------------------------------------------------
    \node[tsnode] (mpa)   at (7.0,1.7)  {MOON\_PA};
    \node[tsnode] (mme)   at (7.0,3.2)  {MOON\_ME};
    \node[tsnode] (mop)   at (9.9,0)    {MOON\_OP};
    % --- planets ------------------------------------------------------
    \node[tsnode] (pci)   at (-0.2,-1.7) [anchor=east] {\footnotesize $\langle$PLANET$\rangle$\_CI};
    \node[tsnode] (pfx)   at (-0.2,-3.2) [anchor=east] {\footnotesize $\langle$PLANET$\rangle$\_FIXED};
    % --- edges --------------------------------------------------------
    \draw[model]    (icrf) -- node[above,font=\scriptsize] {SSB$\to$E} (gcrf);
    \draw[model]    (gcrf) -- node[above,font=\scriptsize] {E$\to$M} (mci);
    \draw[linear]   (gcrf) -- node[left,font=\scriptsize,pos=0.55] {$\mat{R}_{\mathrm{polar}}\mat{R}_{\mathrm{ERA}}\mat{R}_{\mathrm{PN}}$} (itrf);
    \draw[constant] (gcrf) -- node[right,font=\scriptsize,pos=0.55] {$\mat{B}_e$} (eme);
    \draw[model]    (gcrf) -- node[right,font=\scriptsize] {synodic} (emr);
    \draw[linear]   (mci)  -- node[right,font=\scriptsize] {$\mat{R}_{\mathrm{PA},\mathrm{CI}}(\phi,\theta,\psi)$} (mpa);
    \draw[constant] (mpa)  -- node[right,font=\scriptsize] {$\mat{B}_{\mathrm{moon}}$} (mme);
    \draw[linear]   (mci)  -- node[above,font=\scriptsize] {$\mat{R}_{\mathrm{OP}}$} (mop);
    \draw[model]    (icrf) -- (pci);
    \draw[linear]   (pci)  -- node[left,font=\scriptsize] {IAU $W,\alpha,\delta$} (pfx);
    \draw[table]    (gcrf) -- (nyi);
    % --- legend -------------------------------------------------------
    \begin{scope}[shift={(0,-4.6)}]
      \draw[model]    (0,0)   -- (0.8,0)   node[right,font=\scriptsize] {ephemeris translation};
      \draw[linear]   (4.2,0) -- (5.0,0)   node[right,font=\scriptsize] {time-dependent rotation};
      \draw[constant] (8.2,0) -- (9.0,0)   node[right,font=\scriptsize] {constant bias};
      \draw[table]    (4.2,-0.55) -- (5.0,-0.55) node[right,font=\scriptsize] {declared, not implemented};
    \end{scope}
  \end{tikzpicture}
  \caption{The implemented frame graph. ICRF, GCRF and MOON\_CI are the
    three hubs; every conversion is routed as a path through them. Edge
    style encodes the nature of the transformation.}
  \label{fig:fc-graph}
\end{figure}

Solar-system planet frames are peeled off first, because they need the ICRF
hub rather than the GCRF one, and because a same-body \code{<PLANET>\_CI}
$\leftrightarrow$ \code{<PLANET>\_FIXED} pair must \emph{not} round-trip
through the barycenter --- doing so would add and subtract a
$\sim\SI{1e11}{\meter}$ vector and destroy the precision of a small
body-relative state.

\begin{lstlisting}[language=cpp, caption={The routing switch and the
  planet-frame fast path
  (\texttt{cpp/lupnt/conversions/frame\_converter.cc}).},
  label={lst:fc-route}]
// cpp/lupnt/conversions/frame_converter.cc :: ConvertFrameBase
template <int N>
Vec<N> ConvertFrameBase(Real t_tdb, const Vec<N>& rv_in, Frame frame_in, Frame frame_out) {
  if (frame_in == frame_out) return rv_in;
  if (IsPlanetFrame(frame_in) || IsPlanetFrame(frame_out))
    return ConvertPlanetFrame(t_tdb, rv_in, frame_in, frame_out);
  switch (frame_in) {
    case Frame::ICRF:
      return ConvertFrame(t_tdb, IcrfToGcrf(t_tdb, rv_in), Frame::GCRF, frame_out);
    case Frame::ITRF:
      return ConvertFrame(t_tdb, ItrfToGcrf(t_tdb, rv_in), Frame::GCRF, frame_out);
    case Frame::GCRF:
      switch (frame_out) {
        case Frame::ICRF: return GcrfToIcrf(t_tdb, rv_in);
        case Frame::ITRF: return GcrfToItrf(t_tdb, rv_in);
        case Frame::EMR:  return GcrfToEmr(t_tdb, rv_in);
        case Frame::EME:  return GcrfToEme(rv_in);
        case Frame::SER: case Frame::GSE:
        case Frame::MOD: case Frame::TOD: throw std::runtime_error("Not implemented");
        case Frame::MOON_PA: case Frame::MOON_ME: case Frame::MOON_OP: case Frame::MOON_CI:
          return ConvertFrame(t_tdb, GcrfToMoonCi(t_tdb, rv_in), Frame::MOON_CI, frame_out);
        default: throw std::runtime_error("Not implemented");
      }
    case Frame::MOON_CI:
      switch (frame_out) {
        case Frame::MOON_OP: return MoonCiToOp(t_tdb, rv_in);
        case Frame::MOON_PA: return MoonCiToPa(t_tdb, rv_in);
        case Frame::MOON_ME: return MoonPaToMe(MoonCiToPa(t_tdb, rv_in));
        default: return ConvertFrame(t_tdb, MoonCiToGcrf(t_tdb, rv_in), Frame::GCRF, frame_out);
      }
    // ... MOON_ME / MOON_PA / MOON_OP / EME / EMR hop back to their hub ...
  }
}
\end{lstlisting}

\begin{lstlisting}[language=cpp, caption={Same-body fast path for planet
  frames (\texttt{cpp/lupnt/conversions/frame\_converter.cc}).},
  label={lst:fc-planetroute}]
// cpp/lupnt/conversions/frame_converter.cc :: ConvertPlanetFrame
if (IsPlanetFrame(frame_in) && IsPlanetFrame(frame_out)
    && GetFrameCenter(frame_in) == GetFrameCenter(frame_out)) {
  BodyId body = GetFrameCenter(frame_in);
  bool in_fixed = IsPlanetFixedFrame(frame_in);
  bool out_fixed = IsPlanetFixedFrame(frame_out);
  if (in_fixed == out_fixed) return rv_in;            // CI<->CI or FIXED<->FIXED
  if (out_fixed) return BodyCiToFixed(t_tdb, rv_in, body);
  return BodyFixedToCi(t_tdb, rv_in, body);
}
\end{lstlisting}

Because the routing is recursive, a conversion such as ITRF $\to$ MOON\_ME
is realised as the composition
ITRF $\to$ GCRF $\to$ MOON\_CI $\to$ MOON\_PA $\to$ MOON\_ME, i.e.\ four
elementary steps, each contributing its own rotation and, where applicable,
its own $\dot{\mat{R}}\vec{r}$ term.

\section{ICRF and GCRF}
\label{sec:fc-icrf}

GCRF and ICRF share axes and differ only in origin. With
$\vec{r}_{\mathrm{SSB}\to E}$ and $\vec{v}_{\mathrm{SSB}\to E}$ the
barycentric state of the Earth from DE440 \citep{park2021de440}, expressed in
GCRF (= ICRS) axes,
%
\begin{equation}
  \label{eq:fc-icrf}
  \vec{r}_{\mathrm{ICRF}} = \vec{r}_{\mathrm{GCRF}} + \vec{r}_{\mathrm{SSB}\to E},
  \qquad
  \vec{v}_{\mathrm{ICRF}} = \vec{v}_{\mathrm{GCRF}} + \vec{v}_{\mathrm{SSB}\to E},
\end{equation}
%
so $\mat{R}_{\mathrm{ICRF},\mathrm{GCRF}} = \mat{I}$ and
$\vec{p}_{\mathrm{ICRF},\mathrm{GCRF}} = \vec{r}_{\mathrm{SSB}\to E}$; the
inverse subtracts the same state. This is a purely kinematic
(Newtonian) origin shift: no relativistic scale factor between TCB-based and
TCG-based coordinates is applied, which is the correct choice given that
\lupnt{}'s states are TDB-compatible throughout.

\section{GCRF and ITRF}
\label{sec:fc-itrf}

This is the only conversion in \lupnt{} that involves three different time
scales simultaneously. It follows the IAU 2006/2000A CIO-based convention
\citep{iers2010}, factored into the product of three rotations:
%
\begin{equation}
  \label{eq:fc-itrf-chain}
  \mat{R}_{\mathrm{ITRF},\mathrm{GCRF}}(t)
  =
  \underbrace{\mat{R}_{\mathrm{polar}}\!\left(t_{\TT}, x_p, y_p\right)}_{\text{polar motion}}
  \;
  \underbrace{\mat{R}_{\mathrm{ERA}}\!\left(t_{\UTone}\right)}_{\text{Earth rotation}}
  \;
  \underbrace{\mat{R}_{\mathrm{PN}}\!\left(t_{\TT}\right)}_{\text{bias--precession--nutation}} .
\end{equation}
%
Read right to left: a GCRF vector is first carried into the Celestial
Intermediate Reference System, then rotated about the Celestial Intermediate
Pole (CIP) through the Earth Rotation Angle into the Terrestrial
Intermediate Reference System, and finally tilted onto the ITRF pole by
polar motion.

\subsection{Bias, precession and nutation}

\lupnt{} does not evaluate the classical precession and nutation angle
series. It reads the CIP unit-vector coordinates $X, Y$ and the CIO locator
$s$ from a tabulated IAU 2006/2000A (SOFA) series
(\file{cpp/lupnt/interfaces/iau_sofa.cc}, one row per day in arcseconds
against Julian date TT, interpolated with a ninth-order Lagrange
polynomial), and assembles the celestial-to-intermediate matrix in closed
form:
%
\begin{equation}
  \label{eq:fc-pn}
  \mat{R}_{\mathrm{PN}}
  =
  \mat{R}_z(-s)
  \begin{bmatrix}
    1-aX^2 & -aXY   & -X \\
    -aXY   & 1-aY^2 & -Y \\
    X      & Y      & 1-a\left(X^2+Y^2\right)
  \end{bmatrix},
  \qquad
  a = \frac{1}{1+\sqrt{1-X^2-Y^2}} .
\end{equation}
%
Here $X = \sin d\cos E$ and $Y = \sin d\sin E$ are the direction cosines of
the CIP in the GCRS, $d$ and $E$ being its polar coordinates, and $s$ is the
CIO locator that positions the Celestial Intermediate Origin on the
intermediate equator. Because the frame bias (the offset between the GCRS
pole/origin and the mean pole/equinox of J2000) is already folded into the
tabulated $X, Y, s$, \cref{eq:fc-pn} is a \emph{bias--precession--nutation}
matrix: no separate bias factor appears in \cref{eq:fc-itrf-chain}. The
argument is $t_{\TT}$, obtained as
$t_{\TT} = \code{ConvertTime}(t_{\TDB}, \TDB, \TT)$ and converted to a
Julian date.

\subsection{Earth rotation}

$\mat{R}_{\mathrm{ERA}} = \mat{R}_z(\theta_{\mathrm{ERA}})$ is a rotation
about the CIP through the Earth Rotation Angle, which is \emph{exactly}
linear in UT1 by definition \citep{iers2010}:
%
\begin{equation}
  \label{eq:fc-era}
  \theta_{\mathrm{ERA}}(t_{\UTone})
  =
  2\pi\left(0.7790572732640 + 1.00273781191135448\, D_{\UTone}\right),
  \qquad
  D_{\UTone} = \frac{t_{\UTone}}{\SI{86400}{\second}},
\end{equation}
%
with $D_{\UTone}$ the number of UT1 days elapsed since J2000; the result is
wrapped to $(-\pi,\pi]$. UT1 itself is recovered in one of two ways: if a
SPICE-fitted EOP model covers the epoch (\cref{sec:fc-spicefit}), then
$t_{\UTone} = t_{\UTC} + (\UTone-\UTC)_{\text{fitted}}$; otherwise the
generic $\TDB\to\UTone$ conversion of \cref{chap:time-conversions} is used,
which interpolates the bundled IERS EOP~14~C04 table.

$\theta_{\mathrm{ERA}}$ is the only factor of \cref{eq:fc-itrf-chain} whose
derivative is retained. Its rate is the Earth rotation rate
$\omega_\oplus \approx \SI{7.292115e-5}{\radian\per\second}$, computed as
%
\begin{equation}
  \label{eq:fc-omega-earth}
  \omega_\oplus =
  \begin{cases}
    \dfrac{2\pi \times 1.00273781191135448}{\SI{86400}{\second}}
      \left(1 + \dfrac{\dd(\UTone-\UTC)}{\dd t}\right)
      & \text{fitted EOP available,}\\[3.2ex]
    \num{7.292115146706979e-5}\left(1 - \dfrac{\mathrm{LOD}}{\SI{86400}{\second}}\right)
      & \text{otherwise,}
  \end{cases}
\end{equation}
%
where LOD is the excess length of day in seconds interpolated from the EOP
table at UTC. In the fitted branch the derivative of $\UTone-\UTC$ is
obtained analytically by differentiating the Chebyshev polynomial, so it is
exact rather than a finite difference.

\subsection{Polar motion}

$\mat{R}_{\mathrm{polar}}$ carries the CIP onto the ITRF pole using the pole
offsets $x_p, y_p$ (radians, from the fitted EOP or interpolated from the
bundled table at UTC) and the TIO locator $s'$:
%
\begin{equation}
  \label{eq:fc-polar}
  \mat{R}_{\mathrm{polar}}
  = \mat{R}_x(-y_p)\,\mat{R}_y(-x_p)\,\mat{R}_z(s'),
  \qquad
  s' = -47\,\mu''\,T,
  \qquad
  T = \frac{t_{\TT}[\si{\second}]}{86400 \times 36525},
\end{equation}
%
where $T$ is the number of Julian centuries of TT elapsed since J2000, as the
IERS definition of the TIO locator requires \citep[Eq.~(5.13)]{iers2010}.
\lupnt{} carries $t_{\TT}$ as seconds past J2000, so the conversion divides by
\code{DAYS\_CENTURY}~$\times$~\code{SECS\_DAY}. The evaluation is factored into
a single file-local helper, templated on the scalar type so that the
\cpp{Real} caller keeps its autodiff derivative and the \cpp{double} caller
stays plain arithmetic:

\begin{lstlisting}[language=cpp, caption={The one place the TIO locator is
  evaluated (\texttt{cpp/lupnt/conversions/frame\_conversions.cc}).},
  label={lst:fc-tio}]
// cpp/lupnt/conversions/frame_conversions.cc :: TioLocator
template <typename T> inline T TioLocator(const T& t_tt) {
  return -47e-6 * RAD_ARCSEC * (t_tt / (DAYS_CENTURY * SECS_DAY));
}
\end{lstlisting}

\begin{lupntnote}
  Both the division by \code{SECS\_DAY} and the fact that there is only one
  copy of \cref{lst:fc-tio} are corrections. The expression previously divided
  by \code{DAYS\_CENTURY} alone, treating seconds as days: the argument was
  $t_{\TT}[\si{\second}]/36525$ instead of
  $t_{\TT}[\si{\second}]/(86400\times36525)$, inflating $s'$ by a factor of
  $\num{86400}$. By the mid-2020s that gave $s' \approx \num{4.9e-6}$~rad in
  place of the intended $\num{5.7e-11}$~rad --- since $s'$ enters
  $\mat{R}_z$ alongside $\theta_{\mathrm{ERA}}$, about \SI{31}{\meter} of
  spurious rotation about the polar axis at the equator.

  The more useful lesson is why it went unnoticed for so long. The formula was
  written out in \emph{three} places: \cpp{RotPolarMotion}, which applies
  $s'$; \cpp{ComputeEopFromSpice} (\cref{sec:fc-spicefit}), which subtracts it
  again when inverting SPICE for the ERA; and a hardcoded copy inside the
  round-trip test in \file{cpp/test/conversions/test_frame_conversions.cc}. In
  that round trip $s'$ cancels identically --- one site removes exactly what
  the other re-applies --- so the test only ever verified that the copies
  \emph{agreed with each other}. It would have passed for any value of $s'$
  whatsoever, including zero, which is precisely how a wrong one survived.
  The two library sites are now the single helper of \cref{lst:fc-tio}; the
  test keeps its own independent expression on purpose, since a test that
  called the helper would verify nothing. Collapsing the duplicates removes the
  possibility of disagreement but not the possibility of a shared error, so the
  magnitude itself is now pinned separately by
  \code{conversions.frame\_conversions\_extra.tio\_locator\_magnitude} in
  \file{cpp/test/conversions/test_frame_conversions_extra.cc}, which recovers
  $s'$ from $\mat{R}_{\mathrm{polar}}$ and bounds it at
  $\num{1e-8}$~rad --- orders of magnitude away from both the correct value and
  the old one.
\end{lupntnote}

\subsection{The complete mapping and its inverse}

The position and velocity mapping is
%
\begin{align}
  \label{eq:fc-itrf-rv}
  \vec{r}_{\mathrm{ITRF}} &= \mat{R}_{\mathrm{ITRF},\mathrm{GCRF}}\,\vec{r}_{\mathrm{GCRF}}, \\
  \vec{v}_{\mathrm{ITRF}} &= \mat{R}_{\mathrm{ITRF},\mathrm{GCRF}}\,\vec{v}_{\mathrm{GCRF}}
    + \dot{\mat{R}}_{\mathrm{ITRF},\mathrm{GCRF}}\,\vec{r}_{\mathrm{GCRF}},
  \nonumber
\end{align}
%
where only the sidereal factor is differentiated,
%
\begin{equation}
  \label{eq:fc-itrf-rdot}
  \dot{\mat{R}}_{\mathrm{ITRF},\mathrm{GCRF}}
  = \mat{R}_{\mathrm{polar}}\,\dot{\mat{R}}_{\mathrm{ERA}}\,\mat{R}_{\mathrm{PN}} .
\end{equation}
%
Applying \cref{eq:fc-omega} to \cref{eq:fc-itrf-rdot} and using
$\dot{\mat{R}}_{\mathrm{ERA}}\mat{R}_{\mathrm{ERA}}\transpose
= -\left[\omega_\oplus\uvec{z}\right]_\times$ gives
%
\begin{equation}
  \label{eq:fc-itrf-omega}
  \dot{\mat{R}}_{\mathrm{ITRF},\mathrm{GCRF}}\,
  \mat{R}_{\mathrm{ITRF},\mathrm{GCRF}}\transpose
  = -\left[\vec{\omega}_\oplus\right]_\times,
  \qquad
  \vec{\omega}_\oplus = \omega_\oplus\, \mat{R}_{\mathrm{polar}}\uvec{z},
\end{equation}
%
so \cref{eq:fc-itrf-rv} is exactly the classical
$\vec{v}_{\mathrm{ITRF}} = \mat{R}\vec{v}_{\mathrm{GCRF}}
- \vec{\omega}_\oplus\times\vec{r}_{\mathrm{ITRF}}$ with the spin axis taken
as the polar-motion-corrected CIP direction --- to first order, the ITRF
$z$-axis. Neglecting $\dot{\mat{R}}_{\mathrm{PN}}$ and
$\dot{\mat{R}}_{\mathrm{polar}}$ discards precession--nutation rates of
order $\num{1e-12}$~rad/s and polar-motion rates of order $\num{1e-13}$~rad/s,
seven to eight orders of magnitude below $\omega_\oplus$; see
\cref{sec:fc-boundaries} for the resulting velocity error.

The inverse solves \cref{eq:fc-itrf-rv} for the GCRF state, reusing the
already-computed $\vec{r}_{\mathrm{GCRF}}$ in the transport term:
%
\begin{align}
  \label{eq:fc-itrf-inv}
  \vec{r}_{\mathrm{GCRF}} &= \mat{R}_{\mathrm{ITRF},\mathrm{GCRF}}\transpose\,\vec{r}_{\mathrm{ITRF}}, \\
  \vec{v}_{\mathrm{GCRF}} &= \mat{R}_{\mathrm{ITRF},\mathrm{GCRF}}\transpose
    \left(\vec{v}_{\mathrm{ITRF}} - \dot{\mat{R}}_{\mathrm{ITRF},\mathrm{GCRF}}\,\vec{r}_{\mathrm{GCRF}}\right).
  \nonumber
\end{align}

\begin{lstlisting}[language=cpp, caption={GCRF to ITRF, rotation product and
  velocity term
  (\texttt{cpp/lupnt/conversions/frame\_conversions.cc}).},
  label={lst:fc-gcrf-itrf}]
// cpp/lupnt/conversions/frame_conversions.cc :: GcrfToItrf
Vec6 GcrfToItrf(Real t_tdb, const Vec6& rv_gcrf) {
  Mat3 R_po    = RotPolarMotion(t_tdb);        // R_polar(t_TT, x_p, y_p)
  Mat3 R_pn    = RotPrecessionNutation(t_tdb); // R_PN(t_TT) from tabulated (X, Y, s)
  Mat3 R_s     = RotSideralMotion(t_tdb);      // R_z(theta_ERA(t_UT1))
  Mat3 R_s_dot = RotSideralMotionDot(t_tdb);   // d/dt of the above

  Mat3 R_GcrfToItrf     = R_po * R_s     * R_pn;
  Mat3 R_GcrfToItrf_dot = R_po * R_s_dot * R_pn;

  Vec3 r_gcrf = rv_gcrf.head(3);
  Vec3 v_gcrf = rv_gcrf.tail(3);
  Vec3 r_itrf = R_GcrfToItrf * r_gcrf;
  Vec3 v_itrf = R_GcrfToItrf * v_gcrf + R_GcrfToItrf_dot * r_gcrf;

  Vec6 rv_itrf;
  rv_itrf << r_itrf, v_itrf;
  return rv_itrf;
}
\end{lstlisting}

\begin{lstlisting}[language=cpp, caption={Precession--nutation from the
  tabulated CIP coordinates
  (\texttt{cpp/lupnt/conversions/frame\_conversions.cc}).},
  label={lst:fc-pn}]
// cpp/lupnt/conversions/frame_conversions.cc :: RotPrecessionNutation
Real t_tt  = ConvertTime(t_tdb, Time::TDB, Time::TT);
Real jd_tt = TimeToJd(t_tt);
IauSofaData iau_data = GetIauSofaData(jd_tt);   // X, Y, s in arcsec
Real X = iau_data.X * RAD_ARCSEC;
Real Y = iau_data.Y * RAD_ARCSEC;
Real s = iau_data.s * RAD_ARCSEC;
Real a = 1. / (1. + sqrt(1. - X * X - Y * Y));
Mat3 mat{{1. - a * X * X, -a * X * Y,     -X},
         {-a * X * Y,     1. - a * Y * Y, -Y},
         {X,              Y,              1. - a * (X * X + Y * Y)}};
Mat3 R_pn = RotZ(-s) * mat;
\end{lstlisting}

\begin{lupntnote}
  The \cpp{Vec3} (position-only) overloads of \cpp{GcrfToItrf} and
  \cpp{ItrfToGcrf} skip \cpp{RotSideralMotionDot} entirely. They are
  therefore cheaper by one EOP lookup, and are the right entry point for
  static ground-station coordinates.
\end{lupntnote}

\subsection{Equinox-based sidereal angles}

\lupnt{} also provides the classical equinox-based angles in
\file{cpp/lupnt/conversions/time_conversions.cc}, but they are \emph{not}
used by \cpp{GcrfToItrf}. Greenwich Mean Sidereal Time is evaluated from the
IAU-82 polynomial,
%
\begin{equation}
  \label{eq:fc-gmst}
  \begin{aligned}
    \theta_{\mathrm{GMST}}
      &= 2\pi \operatorname{frac}\!\left(\frac{g}{\SI{86400}{\second}}\right), \\
    g &= 24110.54841 + 8640184.812866\,T_0
         + 1.002737909350795\,u
         + \left(0.093104 - \num{6.2e-6}\,T\right)T^2,
  \end{aligned}
\end{equation}
%
with $T_0$ the Julian centuries from J2000 to UT1 midnight of the day in
question, $u$ the seconds of UT1 elapsed since that midnight, and $T$ the
Julian centuries from J2000 to the epoch itself; Greenwich Apparent Sidereal
Time adds the equation of the equinoxes,
%
\begin{equation}
  \label{eq:fc-gast}
  \theta_{\mathrm{GAST}} = \theta_{\mathrm{GMST}} + \Delta\psi\cos\varepsilon
  \pmod{2\pi}.
\end{equation}
%
GAST is used by the Sun/Moon analytic ephemeris helper
(\file{cpp/lupnt/environment/solar_system.cc}, which builds
$\mat{R}_z(\theta_{\mathrm{GAST}})$), and both angles are cross-validated
against Orekit. Because \cref{eq:fc-itrf-chain} is CIO-based, the equinox
route and the CIO route are alternative --- not composable --- descriptions
of the same physical rotation.

\section{GCRF and EME2000}
\label{sec:fc-eme}

EME2000 (also called J2000 or, in the enum, \code{ECI}) is the mean equator
and equinox of J2000. It differs from GCRF by the constant IAU 2006 frame
bias:
%
\begin{equation}
  \label{eq:fc-bias}
  \mat{B}_e \equiv \mat{R}_{\mathrm{EME},\mathrm{GCRF}}
  = \mat{R}_x(-\eta_0)\,\mat{R}_y(\xi_0)\,\mat{R}_z(\Delta\alpha_0),
\end{equation}
%
with the three bias angles taken from
\file{cpp/lupnt/core/constants.h},
%
\begin{equation}
  \label{eq:fc-bias-angles}
  \xi_0 = -16.6170\ \mathrm{mas}, \qquad
  \eta_0 = -6.8192\ \mathrm{mas}, \qquad
  \Delta\alpha_0 = -14.6\ \mathrm{mas},
\end{equation}
%
where $\xi_0, \eta_0$ are the celestial-pole offsets at J2000 and
$\Delta\alpha_0$ is the right-ascension offset of the GCRS origin from the
mean equinox \citep{iers2010}. \Cref{eq:fc-bias} is exact to all orders in
the angles because it is a genuine product of elementary rotations; the
first- and second-order approximations are also available for comparison,
%
\begin{equation}
  \label{eq:fc-bias-first}
  \mat{B}_e \approx
  \begin{bmatrix}
    1 & \Delta\alpha_0 & -\xi_0 \\
    -\Delta\alpha_0 & 1 & -\eta_0 \\
    \xi_0 & \eta_0 & 1
  \end{bmatrix}
  \quad (\cpp{RotGcrfToEmeFirstOrder}),
\end{equation}
%
with \cpp{RotGcrfToEmeSecondOrder} retaining the quadratic terms. Since
$\mat{B}_e$ is constant, $\dot{\mat{B}}_e = \mat{0}$ and the same matrix is
applied to position and velocity; the inverse is
$\mat{B}_e\transpose$. Neither \cpp{GcrfToEme} nor \cpp{EmeToGcrf} takes an
epoch argument.

\section{The Earth--Moon rotating frame}
\label{sec:fc-emr}

\code{EMR} is a synodic frame built on the instantaneous Earth$\to$EMB
(Earth--Moon barycenter) state $(\vec{r}_c, \vec{v}_c)$ taken from DE440 in
GCRF axes. Its basis is
%
\begin{equation}
  \label{eq:fc-emr-basis}
  \uvec{x} = \frac{\vec{r}_c}{\lVert\vec{r}_c\rVert},
  \qquad
  \uvec{n} = \frac{\vec{r}_c\times\vec{v}_c}{\lVert\vec{r}_c\times\vec{v}_c\rVert},
  \qquad
  \uvec{y} = \uvec{n}\times\uvec{x},
\end{equation}
%
and the rotation from GCRF to EMR stacks these as rows in the order
$(\uvec{x}, \uvec{y}, \uvec{n})$, so the EMR components of a vector are
(radial, in-track, orbit-normal). The transform of a GCRF state
$(\vec{r}_d,\vec{v}_d)$ is
%
\begin{equation}
  \label{eq:fc-emr}
  \vec{r}_{\mathrm{EMR}} = \mat{R}\left(\vec{r}_d - \vec{r}_c\right),
  \qquad
  \vec{v}_{\mathrm{EMR}} = \mat{R}\left(\vec{v}_d - \vec{v}_c
    - \vec{\omega}\times\left(\vec{r}_d - \vec{r}_c\right)\right),
\end{equation}
%
implemented by \cpp{InertialToSynodic} /\cpp{SynodicToInertial} in
\file{cpp/lupnt/conversions/state_conversions.cc}, which \cpp{GcrfToEmr} and
\cpp{EmrToGcrf} wrap. The origin of EMR is therefore the EMB, not the Earth.

\begin{lupntwarning}
  \cpp{InertialToSynodic} forms the frame angular velocity as
  $\vec{\omega} = (\vec{r}_c\times\vec{v}_c)/\lVert\vec{r}_c\rVert$, whereas
  the angular velocity of the rotating frame is
  $\vec{\omega} = (\vec{r}_c\times\vec{v}_c)/\lVert\vec{r}_c\rVert^2$. The
  transport term in \cref{eq:fc-emr} is consequently scaled by
  $\lVert\vec{r}_c\rVert$. EMR \emph{positions} are unaffected (they do not
  involve $\vec{\omega}$); EMR \emph{velocities} should not be relied on
  until this is reconciled. The \cpp{Vec3} overloads zero the velocity
  before calling through, so they return correct positions.
\end{lupntwarning}

\section{GCRF and MOON\_CI}
\label{sec:fc-moonci}

\code{MOON\_CI} is Moon-centered and axis-aligned with GCRF/ICRF: the
conversion is a pure translation by the geocentric Moon state, evaluated
from the DE440 Chebyshev coefficients \citep{park2021de440} bundled with
\lupnt{} (\file{cpp/lupnt/interfaces/kernels.cc}). With
$\vec{r}_{E\to M}, \vec{v}_{E\to M}$ the Moon state relative to Earth in
GCRF axes,
%
\begin{equation}
  \label{eq:fc-moonci}
  \vec{r}_{\mathrm{MOON\_CI}} = \vec{r}_{\mathrm{GCRF}} - \vec{r}_{E\to M},
  \qquad
  \vec{v}_{\mathrm{MOON\_CI}} = \vec{v}_{\mathrm{GCRF}} - \vec{v}_{E\to M},
\end{equation}
%
so $\mat{R}_{\mathrm{MOON\_CI},\mathrm{GCRF}} = \mat{I}$ and
$\vec{p} = -\vec{r}_{E\to M}$. The velocity is the analytic derivative of
the position Chebyshev series, not a finite difference.

\begin{lstlisting}[language=cpp, caption={GCRF to Moon-centered inertial ---
  a pure DE440 translation
  (\texttt{cpp/lupnt/conversions/frame\_conversions.cc}).},
  label={lst:fc-moonci}]
// cpp/lupnt/conversions/frame_conversions.cc :: GcrfToMoonCi
Vec6 GcrfToMoonCi(Real t_tdb, const Vec6& rv_gcrf) {
  Vec6 rv_earth2moon = GetBodyPosVel(t_tdb, BodyId::EARTH, BodyId::MOON, Frame::GCRF);
  Vec6 rv_mi = rv_gcrf - rv_earth2moon;
  return rv_mi;
}
\end{lstlisting}

\section{MOON\_CI and MOON\_PA}
\label{sec:fc-moonpa}

\code{MOON\_PA} is the lunar principal-axis (mantle) frame. Its orientation
relative to the inertial axes is given by the three Euler angles of the
lunar physical libration, $\phi$ (node), $\theta$ (inclination of the
principal-axis equator) and $\psi$ (rotation), which DE440 supplies as a
Chebyshev series in the same file as the planetary positions
(\cpp{GetLunarMantleData}, ephemeris item
\code{MOON\_MANTLE\_LIBRATIONS}). The rotation is the classical
$Z$--$X$--$Z$ sequence
%
\begin{equation}
  \label{eq:fc-moonpa}
  \mat{R}_{\mathrm{PA},\mathrm{CI}}(t)
  = \mat{R}_z\!\left(\psi(t)\right)\,
    \mat{R}_x\!\left(\theta(t)\right)\,
    \mat{R}_z\!\left(\phi(t)\right),
\end{equation}
%
whose exact derivative follows from the product rule and the elementary
derivative matrices of \cref{eq:fc-elementary}:
%
\begin{equation}
  \label{eq:fc-moonpa-dot}
  \dot{\mat{R}}_{\mathrm{PA},\mathrm{CI}}
  = \mat{R}_z'(\psi,\dot\psi)\,\mat{R}_x(\theta)\,\mat{R}_z(\phi)
  + \mat{R}_z(\psi)\,\mat{R}_x'(\theta,\dot\theta)\,\mat{R}_z(\phi)
  + \mat{R}_z(\psi)\,\mat{R}_x(\theta)\,\mat{R}_z'(\phi,\dot\phi).
\end{equation}
%
The angle rates $\dot\phi,\dot\theta,\dot\psi$ come from differentiating the
same Chebyshev series, so \cref{eq:fc-moonpa-dot} is exact (subject to the
ephemeris fit), not a numerical derivative. Unlike the Earth case, no term
is dropped: the lunar rotation rate
($\approx\SI{2.66e-6}{\radian\per\second}$) and the libration rates are of
comparable smallness, so retaining all three products matters. The state
transform and its inverse are
%
\begin{align}
  \label{eq:fc-moonpa-vel}
  \vec{r}_{\mathrm{PA}} &= \mat{R}_{\mathrm{PA},\mathrm{CI}}\,\vec{r}_{\mathrm{CI}},
  &\vec{v}_{\mathrm{PA}} &= \mat{R}_{\mathrm{PA},\mathrm{CI}}\,\vec{v}_{\mathrm{CI}}
    + \dot{\mat{R}}_{\mathrm{PA},\mathrm{CI}}\,\vec{r}_{\mathrm{CI}}, \\
  \vec{r}_{\mathrm{CI}} &= \mat{R}_{\mathrm{PA},\mathrm{CI}}\transpose\,\vec{r}_{\mathrm{PA}},
  &\vec{v}_{\mathrm{CI}} &= \mat{R}_{\mathrm{PA},\mathrm{CI}}\transpose
    \left(\vec{v}_{\mathrm{PA}} - \dot{\mat{R}}_{\mathrm{PA},\mathrm{CI}}\,\vec{r}_{\mathrm{CI}}\right).
  \nonumber
\end{align}

\begin{lstlisting}[language=cpp, caption={Lunar libration rotation and its
  analytic derivative
  (\texttt{cpp/lupnt/conversions/frame\_conversions.cc}).},
  label={lst:fc-moonpa}]
// cpp/lupnt/conversions/frame_conversions.cc :: RotMoonCiToPa
Vec6 lunar_mantle = GetLunarMantleData(t_tdb, compute_vel);
auto [phi, theta, psi, phi_dot, theta_dot, psi_dot] = Unpack(lunar_mantle);

Mat3 R_z_phi = RotZ(phi), R_x_theta = RotX(theta), R_z_psi = RotZ(psi);
Mat3 R_mi2pa = R_z_psi * R_x_theta * R_z_phi;

if (compute_vel) {
  *R_mi2pa_dot = RotZdot(psi, psi_dot) * R_x_theta          * R_z_phi
               + R_z_psi               * RotXdot(theta, theta_dot) * R_z_phi
               + R_z_psi               * R_x_theta          * RotZdot(phi, phi_dot);
}
\end{lstlisting}

Two implementation details are worth recording. First,
\cpp{GetLunarMantleData} divides the raw ephemeris item by
\code{M\_KM}~$=\num{1e3}$, undoing the kilometre-to-metre scaling that the
generic Chebyshev reader applies to \emph{position} items, so the returned
angles are radians and the rates are \si{\radian\per\second}. Second,
\cpp{RotMoonCiToPa} caches the last evaluated $(\mat{R},\dot{\mat{R}})$ pair
keyed on the epoch, because a single \cpp{ConvertFrame} call from a
lunar-fixed frame to an Earth frame can request it several times.

\begin{lupntnote}
  The \code{M\_KM} division above is applied \emph{once}, in the scalar
  \cpp{GetLunarMantleData(Real, bool)} overload. The vectorised overload
  \cpp{GetLunarMantleData(VecX, bool)} in
  \file{cpp/lupnt/interfaces/kernels.cc} used to divide the scalar overload's
  return value by \code{M\_KM} a second time, so every libration angle it
  produced was $\num{1e3}$ too small: at one sample epoch it returned
  $\phi = \num{-5.41e-5}$ where the scalar overload returned
  $\num{-5.41e-2}$. The defect was latent --- nothing inside \lupnt{} routes
  \cref{eq:fc-moonpa} through the vector overload, so no \code{MOON\_PA}
  result was ever wrong --- but the overload is declared public API in
  \file{cpp/lupnt/interfaces/kernels.h}, so any external caller batching
  epochs would have received silently wrong angles. The
  extra division has been removed, and the two overloads are now required to
  agree row for row by
  \code{interfaces.kernels\_extra.lunar\_mantle\_vector\_matches\_scalar} in
  \file{cpp/test/interfaces/test_kernels_extra.cc}.
\end{lupntnote}

\begin{lupntnote}
  If \cpp{InitFrameConversionFromSpice} has fitted a lunar-orientation model
  covering $t_{\TDB}$, \cpp{TryGetFittedMoonCiToPa} returns the
  SPICE-derived rotation and its exact derivative instead, and the DE
  libration branch is bypassed entirely (\cref{sec:fc-spicefit}).
\end{lupntnote}

\section{MOON\_PA and MOON\_ME}
\label{sec:fc-moonme}

The mean-Earth/polar-axis frame \code{MOON\_ME} --- the frame in which lunar
cartographic products and surface-station coordinates are published --- is
related to the principal-axis frame by a constant bias rotation of three
small angles:
%
\begin{equation}
  \label{eq:fc-moonme}
  \mat{B}_{\mathrm{moon}} \equiv \mat{R}_{\mathrm{ME},\mathrm{PA}}
  = \mat{R}_x(-0.2785'')\,\mat{R}_y(-78.6944'')\,\mat{R}_z(-67.8526'') .
\end{equation}
%
The total angle is about $104''$, i.e.\ roughly \SI{880}{\meter} on the
lunar surface --- large enough that confusing PA and ME coordinates is a
first-order error for any surface user. As with the Earth frame bias, the
same matrix is applied to position and velocity, no epoch is needed, and the
inverse is the transpose.

\section{MOON\_OP: the lunar orbit-plane frame}
\label{sec:fc-moonop}

\code{MOON\_OP} is the frozen-orbit analysis frame: it is Moon-centered,
with its $z$-axis along the Earth--Moon orbital angular momentum and its
$x$-axis the lunar mean-Earth pole projected into that plane. Let
$(\vec{r}_{M\to E}, \vec{v}_{M\to E})$ be the Moon-to-Earth state in GCRF
axes; then
%
\begin{equation}
  \label{eq:fc-moonop-z}
  \uvec{z}_{\mathrm{OP}}
  = \frac{\vec{r}_{M\to E}\times\vec{v}_{M\to E}}
         {\lVert\vec{r}_{M\to E}\times\vec{v}_{M\to E}\rVert},
\end{equation}
%
and with $\uvec{p}$ the \code{MOON\_ME} $z$-axis expressed in
\code{MOON\_CI} (obtained by a recursive \cpp{ConvertFrame} call on
$\uvec{e}_z$),
%
\begin{equation}
  \label{eq:fc-moonop-xy}
  \uvec{x}_{\mathrm{OP}}
  = \frac{\uvec{p}\times\uvec{z}_{\mathrm{OP}}}
         {\lVert\uvec{p}\times\uvec{z}_{\mathrm{OP}}\rVert},
  \qquad
  \uvec{y}_{\mathrm{OP}} = \uvec{z}_{\mathrm{OP}}\times\uvec{x}_{\mathrm{OP}} .
\end{equation}
%
The columns of the OP-to-CI rotation are these three unit vectors expressed
in CI axes,
%
\begin{equation}
  \label{eq:fc-moonop-R}
  \mat{R}_{\mathrm{CI},\mathrm{OP}}
  = \begin{bmatrix}
      \uvec{x}_{\mathrm{OP}} & \uvec{y}_{\mathrm{OP}} & \uvec{z}_{\mathrm{OP}}
    \end{bmatrix},
\end{equation}
%
and \cpp{MoonCiToOp} applies its transpose.

\begin{lstlisting}[language=cpp, caption={Orbit-plane triad
  (\texttt{cpp/lupnt/conversions/frame\_conversions.cc}).},
  label={lst:fc-moonop}]
// cpp/lupnt/conversions/frame_conversions.cc :: RotMoonOpToCi
Vec6 rv_m2e = GetBodyPosVel(t_tdb, BodyId::MOON, BodyId::EARTH, Frame::GCRF);
Vec3 z_op   = rv_m2e.head(3).cross(rv_m2e.tail(3)).normalized();
Vec3 i_pole = ConvertFrame(t_tdb, Vec3::UnitZ(), Frame::MOON_ME, Frame::MOON_CI);
Vec3 x_op   = i_pole.cross(z_op).normalized();
Vec3 y_op   = z_op.cross(x_op).normalized();
Mat3 R_op2ci;
R_op2ci << x_op, y_op, z_op;
\end{lstlisting}

\begin{lupntwarning}
  \cpp{MoonCiToOp}/\cpp{MoonOpToCi} apply $\mat{R}_{\mathrm{CI},\mathrm{OP}}$
  to position \emph{and} velocity with no $\dot{\mat{R}}$ term. The OP frame
  precesses with the lunar orbit at
  $\sim\SI{2.7e-6}{\radian\per\second}$, so an OP velocity omits a transport
  term of order $\omega r$ --- about \SI{5}{\meter\per\second} at
  \SI{2000}{\kilo\meter} selenocentric radius. \code{MOON\_OP} is intended
  for geometric orbit-design work (orbit-plane orientation, frozen-orbit
  elements), not for velocity-critical computation.
\end{lupntwarning}

\section{Solar-system planet frames}
\label{sec:fc-planets}

Each \code{<PLANET>\_CI} frame is planet-centered with ICRF axes; reaching it
from the ICRF hub is a pure ephemeris translation
(\cpp{PlanetCiToIcrf}/\cpp{IcrfToPlanetCi}). Each \code{<PLANET>\_FIXED}
frame co-rotates with the planet under the IAU linear orientation model:
with pole right ascension $\alpha$, pole declination $\delta$, and
prime-meridian angle $W$,
%
\begin{equation}
  \label{eq:fc-iau}
  \mat{R}_{\mathrm{FIXED},\mathrm{CI}}
  = \mat{R}_z(W)\,
    \mat{R}_x\!\left(\tfrac{\pi}{2}-\delta\right)\,
    \mat{R}_z\!\left(\tfrac{\pi}{2}+\alpha\right),
\end{equation}
%
where
%
\begin{equation}
  \label{eq:fc-iau-elements}
  \alpha = \alpha_0 + \dot{\alpha}\,T, \qquad
  \delta = \delta_0 + \dot{\delta}\,T, \qquad
  W = W_0 + \dot{W}\,d,
\end{equation}
%
with $d$ the days and $T = d/36525$ the Julian centuries past J2000 in TT.
The coefficients (\cref{tab:fc-iau}) are the IAU Working Group on
Cartographic Coordinates and Rotational Elements values; the small periodic
terms of the full model are omitted.

\begin{table}[H]
  \centering
  \caption{IAU linear orientation coefficients used by
    \cpp{RotBodyCiToFixed}. Angles in degrees; $T$ in Julian centuries and
    $d$ in days past J2000 (TT).}
  \label{tab:fc-iau}
  \begin{tabular}{l r r r r r r}
    \toprule
    Body & $\alpha_0$ & $\dot{\alpha}$ & $\delta_0$ & $\dot{\delta}$ & $W_0$ & $\dot{W}$ \\
    \midrule
    Mercury & $281.0103$  & $-0.0328$   & $61.4155$  & $-0.0049$ & $329.5988$ & $6.1385108$    \\
    Venus   & $272.76$    & $0$         & $67.16$    & $0$       & $160.20$   & $-1.4813688$   \\
    Mars    & $317.68143$ & $-0.1061$   & $52.88650$ & $-0.0609$ & $176.630$  & $350.89198226$ \\
    Jupiter & $268.056595$& $-0.006499$ & $64.495303$& $0.002413$& $284.95$   & $870.5360000$  \\
    Saturn  & $40.589$    & $-0.036$    & $83.537$   & $-0.004$  & $38.90$    & $810.7939024$  \\
    Uranus  & $257.311$   & $0$         & $-15.175$  & $0$       & $203.81$   & $-501.1600928$ \\
    Neptune & $299.36$    & $0$         & $43.46$    & $0$       & $249.978$  & $541.1397757$  \\
    \bottomrule
  \end{tabular}
\end{table}

Only the prime-meridian spin contributes appreciably to the derivative --- the
pole rates are of order $\SI{0.01}{\degree}$ per century, i.e.\
$\sim\num{1e-16}$~rad/s --- so
%
\begin{equation}
  \label{eq:fc-iau-dot}
  \dot{\mat{R}}_{\mathrm{FIXED},\mathrm{CI}}
  \approx \mat{R}_z'\!\left(W, \dot{W}\right)
    \mat{R}_x\!\left(\tfrac{\pi}{2}-\delta\right)
    \mat{R}_z\!\left(\tfrac{\pi}{2}+\alpha\right),
  \qquad
  \dot{W}\ \text{in \si{\radian\per\second}}.
\end{equation}

\begin{lstlisting}[language=cpp, caption={IAU body orientation
  (\texttt{cpp/lupnt/conversions/frame\_conversions.cc}).},
  label={lst:fc-iau}]
// cpp/lupnt/conversions/frame_conversions.cc :: RotBodyCiToFixed
Real t_tt = ConvertTime(t_tdb, Time::TDB, Time::TT);
Real d = TimeToJd(t_tt) - JD_J2000_TT;  // days past J2000 (TT)
Real T = d / 36525.0;                   // Julian centuries past J2000 (TT)

Real ra  = (o.ra0  + o.ra0_T  * T) * RAD;
Real dec = (o.dec0 + o.dec0_T * T) * RAD;
Real W   = (o.w0   + o.w_dot  * d) * RAD;

Mat3 Rx_pole = RotX(M_PI / 2 - dec);
Mat3 Rz_node = RotZ(M_PI / 2 + ra);
Mat3 R = RotZ(W) * Rx_pole * Rz_node;
if (R_dot != nullptr) {
  Real w_dot = o.w_dot * RAD / SECS_DAY;  // [rad/s]
  *R_dot = RotZdot(W, w_dot) * Rx_pole * Rz_node;
}
\end{lstlisting}

\section{Topocentric frames: ENU and azimuth--elevation--range}
\label{sec:fc-topocentric}

The topocentric helpers in
\file{cpp/lupnt/conversions/coordinate_conversions.cc} operate
\emph{entirely inside one body-fixed frame}; they take a reference position
$\vec{r}_{\mathrm{ref}}$ and a body shape, never a \cpp{Frame} argument. The
caller is responsible for having converted both the target and the reference
into the same body-fixed frame (ITRF for Earth, MOON\_PA or MOON\_ME for the
Moon) with \cpp{ConvertFrame} first.

\paragraph{Geodetic coordinates.} With equatorial radius $R_b$, flattening
$f$, and first eccentricity squared $e^2 = f(2-f)$, the geodetic-to-Cartesian
map is
%
\begin{equation}
  \label{eq:fc-lla}
  \vec{r} =
  \begin{bmatrix}
    (N+h)\cos\varphi\cos\lambda \\
    (N+h)\cos\varphi\sin\lambda \\
    \left((1-e^2)N+h\right)\sin\varphi
  \end{bmatrix},
  \qquad
  N = \frac{R_b}{\sqrt{1-e^2\sin^2\varphi}},
\end{equation}
%
$\varphi$ being geodetic latitude, $\lambda$ longitude and $h$ altitude
above the ellipsoid. The inverse (\cpp{CartToLatLonAlt}) is the classical
fixed-point iteration on $\Delta z = N e^2 \sin\varphi$, iterated to
\code{EPS} and capped at 1000 iterations; passing $R_b = 0$ makes it fall
back to a sphere of radius $\lVert\vec{r}\rVert$.

\paragraph{East--north--up.} The rotation from body-fixed axes to the local
ENU triad at $(\varphi,\lambda)$ is
%
\begin{equation}
  \label{eq:fc-enu-rot}
  \mat{R}_{\mathrm{ENU}}
  = \mat{R}_x\!\left(\tfrac{\pi}{2}-\varphi\right)
    \mat{R}_z\!\left(\tfrac{\pi}{2}+\lambda\right)
  = \begin{bmatrix}
      -\sin\lambda & \cos\lambda & 0 \\
      -\sin\varphi\cos\lambda & -\sin\varphi\sin\lambda & \cos\varphi \\
      \cos\varphi\cos\lambda & \cos\varphi\sin\lambda & \sin\varphi
    \end{bmatrix},
\end{equation}
%
whose rows are the east, north and up unit vectors. The transform of a
state is
%
\begin{equation}
  \label{eq:fc-enu}
  \vec{r}_{\mathrm{ENU}} = \mat{R}_{\mathrm{ENU}}\left(\vec{r}-\vec{r}_{\mathrm{ref}}\right),
  \qquad
  \vec{v}_{\mathrm{ENU}} = \mat{R}_{\mathrm{ENU}}\,\vec{v} .
\end{equation}
%
No $\dot{\mat{R}}$ term appears, and none is needed: $\mat{R}_{\mathrm{ENU}}$
is constant \emph{in the body-fixed frame}, and the input velocity is
already body-fixed-relative. The rotation of the body has been accounted for
once, in the GCRF$\to$ITRF (or MOON\_CI$\to$MOON\_PA) step.

\paragraph{Azimuth, elevation, range.} Finally,
%
\begin{equation}
  \label{eq:fc-azel}
  \rho = \lVert\vec{r}_{\mathrm{ENU}}\rVert, \qquad
  A = \atantwo(e, n), \qquad
  E = \arcsin\!\left(\frac{u}{\rho}\right),
\end{equation}
%
so azimuth is measured clockwise from north and elevation upward from the
local horizon; the inverse is
$(e,n,u) = \rho(\cos E\sin A,\ \cos E\cos A,\ \sin E)$.
\cpp{CartToAzElRange} is the composition of \cref{eq:fc-enu} and
\cref{eq:fc-azel}.

\section{SPICE-calibrated Earth and lunar orientation}
\label{sec:fc-spicefit}

\lupnt{}'s native Earth and lunar orientation models are analytic and
self-contained: an IAU 2006/2000A series plus an IERS EOP table for the
Earth, and DE Chebyshev libration angles for the Moon. Both can disagree
with SPICE's high-accuracy binary PCK products
(\file{earth_latest_high_prec.bpc}, \file{moon_pa_de*.bpc})
\citep{acton1996spice}, most severely when the bundled EOP table does not
cover the requested epoch. \cpp{InitFrameConversionFromSpice} closes that
gap for a known simulation window, without requiring SPICE at evaluation
time.

For the Earth, \cpp{ComputeEopFromSpice} inverts the decomposition
\cref{eq:fc-itrf-chain} at a single epoch. Given SPICE's
\code{"J2000"}$\to$\code{"ITRF93"} rotation $\mat{R}_{\mathrm{SPICE}}$ and
\lupnt{}'s own $\mat{R}_{\mathrm{PN}}$, it forms
$\mat{M} = \mat{R}_{\mathrm{SPICE}}\mat{R}_{\mathrm{PN}}\transpose$ and reads
off
%
\begin{equation}
  \label{eq:fc-eopinv}
  x_p = \arcsin M_{13}, \qquad
  y_p = -\atantwo\!\left(M_{23}, M_{33}\right), \qquad
  s' + \theta_{\mathrm{ERA}} = \atantwo\!\left(M_{12}, M_{11}\right),
\end{equation}
%
then linearises \cref{eq:fc-era} about $t_{\UTone}\approx t_{\TDB}$ to
recover $\UTone-\UTC$ (valid because the true offset is at most about
\SI{1}{\second}, far shorter than the one-day period of
$\theta_{\mathrm{ERA}}$).

Both the three Earth EOP parameters $(x_p, y_p, \UTone-\UTC)$ and the three
lunar Euler angles are then fitted over the requested window with piecewise
Chebyshev polynomials sampled at Chebyshev--Gauss nodes,
%
\begin{equation}
  \label{eq:fc-cheb}
  q_k(t) \approx \sum_{n=0}^{N-1} c_{kn} T_n(\tau),
  \qquad
  \tau = \frac{t - t_{\mathrm{mid}}}{\Delta t/2} \in [-1,1],
\end{equation}
%
with $N = 13$ coefficients per quantity per segment and a default segment
length of \SI{1}{\dayunit}. For the lunar model the fitted quantities are
the ZXZ angles $(\phi,\theta,\psi)$ extracted from SPICE's
\code{"J2000"}$\to$\code{"IAU\_MOON"} rotation by
%
\begin{equation}
  \label{eq:fc-zxz}
  \theta = \arccos R_{33}, \qquad
  \phi = \atantwo\!\left(R_{31}, -R_{32}\right), \qquad
  \psi = \atantwo\!\left(R_{13}, R_{23}\right),
\end{equation}
%
with $\phi$ and $\psi$ continuously unwrapped across the window so that the
polynomial sees a smooth function rather than a $2\pi$ sawtooth. At
evaluation time the rotation is \emph{reconstructed} as
$\mat{R} = \mat{R}_z(\psi)\mat{R}_x(\theta)\mat{R}_z(\phi)$ and its
derivative by the product rule with the analytically differentiated
Chebyshev coefficients. This is the reason the angles, and not the nine
matrix entries, are fitted: a product of elementary rotations is exactly
orthonormal to floating-point precision regardless of the fit residual,
whereas fitting the entries independently left an orthonormality violation
of order $\num{4e-8}$.

\begin{lstlisting}[language=cpp, caption={Fitting the SPICE-derived Earth
  and lunar orientation
  (\texttt{cpp/lupnt/conversions/frame\_conversions.cc}).},
  label={lst:fc-spicefit}]
// cpp/lupnt/conversions/frame_conversions.cc :: InitFrameConversionFromSpice
ChebyshevFitModel earth_fit = FitChebyshevModel(
    [](double t) -> VecXd {
      SpiceEopParams p = ComputeEopFromSpice(Real(t));
      VecXd f(3);
      f << p.x_pole.val(), p.y_pole.val(), p.ut1_utc.val();
      return f;
    },
    t0, t1, /*num_dims=*/3, seg_len, num_coeffs);

ChebyshevFitModel moon_fit = FitEulerAngleModel(
    [](double t) { return SampleSpiceRotationMatrix(t, "J2000", "IAU_MOON"); },
    t0, t1, seg_len, num_coeffs);
\end{lstlisting}

The feature is purely additive: outside the fitted window, or if
\cpp{InitFrameConversionFromSpice} was never called,
\cpp{RotPolarMotion}/\cpp{RotSideralMotion}/\cpp{RotSideralMotionDot} and
\cpp{RotMoonCiToPa} behave exactly as specified in
\cref{sec:fc-itrf,sec:fc-moonpa}. \cpp{ClearFrameConversionFit} removes the
fit; \brk{HasFittedEarthOrientation} and \brk{HasFittedLunarOrientation}
report coverage. A separate, always-SPICE path
(\longid{lupnt::spice::ConvertFrameSpice} in
\file{cpp/lupnt/conversions/frame_converter_spice.cc}) exists as the
reference against which the native path is validated.

\section{Coordinate scale}
\label{sec:fc-units}

Frame conversions do not rescale. Non-dimensionalisation, when used, happens
in the ephemeris layer: for a constant coordinate scale factor $F$,
%
\begin{equation}
  \label{eq:fc-scale}
  \vec{r}_{\mathrm{to}} = F\,\vec{r}_{\mathrm{from}},
  \qquad
  \mu_{\mathrm{to}} = F\,\mu_{\mathrm{from}},
  \qquad
  \vec{v}_{\mathrm{to}} = \vec{v}_{\mathrm{from}} ,
\end{equation}
%
i.e.\ lengths and gravitational parameters scale together while velocities
are invariant. Callers must keep the scale consistent across the states they
transform; see the caveat at the end of \cref{sec:fc-contract}.

\section{The Python binding surface}
\label{sec:fc-python}

\pylupnt{} exposes the \cpp{Frame} enum and the dispatchers directly; every
\cpp{ConvertFrame} overload maps to the single overloaded name
\py{convert\_frame}, with the epoch first and \py{rotate\_only} keyword-only
by convention.

\begin{lstlisting}[language=cpp, caption={Binding definitions
  (\texttt{python/bindings/py\_frame\_converter.cc}).},
  label={lst:fc-bindings}]
// python/bindings/py_frame_converter.cc
py::enum_<Frame>(m, "Frame", "Reference frames (Earth, Moon, and solar-system planet frames).")
    .value("ITRF", Frame::ITRF, "International Terrestrial Reference Frame (Earth-fixed).")
    .value("GCRF", Frame::GCRF, "Geocentric Celestial Reference Frame (Earth-centered inertial).")
    .value("MOON_CI", Frame::MOON_CI, "Moon-centered inertial frame (ICRF-aligned axes).")
    .value("MOON_PA", Frame::MOON_PA, "Moon-fixed principal-axis frame.")
    .value("MOON_ME", Frame::MOON_ME, "Moon-fixed mean-Earth/polar-axis frame.")
    /* ... EME/ECI, ICRF, EMR, MOON_OP, <PLANET>_CI, <PLANET>_FIXED ... */
    .export_values();

m.def("convert_frame", py::overload_cast<Real, const Vec6&, Frame, Frame, bool>(&ConvertFrame),
      py::arg("t_tdb"), py::arg("rv_in"), py::arg("frame_in"), py::arg("frame_out"),
      py::arg("rotate_only") = false);
m.def("get_frame_rotation_translation",
      py::overload_cast<Real, Frame, Frame>(&GetFrameRotationTranslation),
      py::arg("t_tdb"), py::arg("frame_in"), py::arg("frame_out"));
m.def("get_frame_rotation_translation_rv",
      py::overload_cast<Real, Frame, Frame>(&GetFrameRotationTranslationRv),
      py::arg("t_tdb"), py::arg("frame_in"), py::arg("frame_out"));
m.def("get_frame_center", &GetFrameCenter, py::arg("frame"));
m.def("compute_eop_from_spice", /* ... */);      // (x_pole, y_pole, UT1-UTC) at t_tdb
m.def("get_lunar_orientation_angles", /* ... */); // (phi, theta, psi, and rates) at t_tdb
\end{lstlisting}

Eight \py{convert\_frame} overloads are registered, covering the four
argument shapes (single \code{Vec3}/\code{Vec6}, row-wise
\code{MatX3}/\code{MatX6}) against the two epoch shapes (a scalar epoch
applied to all rows, and a vector of epochs applied row by row). The
vectorised forms are the ones to use from Python: they keep the whole loop
inside C++ and, for a fixed epoch, evaluate the ephemeris and the
orientation rotations once.

\begin{lstlisting}[language=py, caption={Typical use from Python.},
  label={lst:fc-pyuse}]
import numpy as np
import pylupnt as pnt

t_tdb = pnt.gregorian_to_time(2026, 7, 21, 0, 0, 0.0)   # [s, TDB past J2000]
rv_gcrf = np.array([6.8e6, 0.0, 0.0, 0.0, 7.6e3, 0.0])  # [m, m/s]

rv_itrf = pnt.convert_frame(t_tdb, rv_gcrf, pnt.Frame.GCRF, pnt.Frame.ITRF)
rv_pa   = pnt.convert_frame(t_tdb, rv_gcrf, pnt.Frame.GCRF, pnt.Frame.MOON_PA)

R, p = pnt.get_frame_rotation_translation(t_tdb, pnt.Frame.GCRF, pnt.Frame.MOON_ME)
assert np.allclose(R @ R.T, np.eye(3), atol=1e-12)
\end{lstlisting}

\section{Model boundaries}
\label{sec:fc-boundaries}

\subsection{What is approximated}

\begin{description}[leftmargin=0.9cm,style=nextline,font=\normalfont\itshape]
  \item[Velocity of the GCRF$\leftrightarrow$ITRF transform.]
    Only $\dot{\mat{R}}_{\mathrm{ERA}}$ is retained
    (\cref{eq:fc-itrf-rdot}). Precession--nutation rates
    ($\sim\num{1e-12}$~rad/s) and polar-motion rates ($\sim\num{1e-13}$~rad/s)
    are dropped, contributing an error below
    $\SI{1e-5}{\meter\per\second}$ at geostationary radius and far less in
    LEO.
  \item[TIO locator.] $s'$ follows the linear IERS model
    $-47\,\mu''\,T$ exactly (\cref{eq:fc-polar}); at $\sim\num{6e-11}$~rad it
    is negligible against every other term in \cref{tab:fc-accuracy}. The
    seconds-versus-centuries unit error that once inflated it by
    $\num{86400}$ has been fixed and is pinned by a magnitude test; see the
    note in \cref{sec:fc-itrf}.
  \item[MOON\_OP velocity.] No $\dot{\mat{R}}$ term
    (\cref{sec:fc-moonop}); errors of order
    $\omega_{\mathrm{Moon}} r \approx \SI{5}{\meter\per\second}$ at
    \SI{2000}{\kilo\meter}.
  \item[EMR velocity.] The synodic angular velocity is formed with
    $\lVert\vec{r}_c\rVert$ instead of $\lVert\vec{r}_c\rVert^2$
    (\cref{sec:fc-emr}). Positions are correct.
  \item[Earth rotation rate.] The two branches of
    \cref{eq:fc-omega-earth} use nominal rates that differ by
    $\approx\num{7e-13}$~rad/s (a relative $\num{1e-8}$), i.e.\
    $\sim\SI{3e-5}{\meter\per\second}$ at GEO.
  \item[Planet orientation.] The IAU model is linear only; the periodic
    terms (Mars, Neptune, and the outer-planet nutations) are omitted. The
    residual is well under a degree over multi-decade spans, which is
    adequate for geometry but not for planetary geodesy.
  \item[EOP interpolation across leap seconds.] The bundled EOP table stores
    daily $\UTone-\UTC$ and interpolates linearly \emph{across} the
    \SI{1}{\second} jump. Within about a day of an insertion this produces a
    spurious offset --- measured at $+\SI{0.938}{\second}$ one hour before the
    2016-12-31 leap and $+\SI{0.094}{\second}$ six hours after the
    2012-06-30 leap --- which propagates directly into
    $\theta_{\mathrm{ERA}}$ and hence into every ITRF result.
  \item[Ephemeris interpolation.] The DE440 Chebyshev evaluation is exact to
    the distributed coefficients; the residual against an independent
    DE440 reader is $\sim\num{5e-11}$ relative.
\end{description}

\subsection{What is not implemented}

\code{MOD}, \code{TOD}, \code{SER} and \code{GSE} are enum values with no
conversion; requesting them throws. \code{EMR} has no
\cpp{frame\_centers} entry. There is no relativistic scale transformation
between TCB-, TDB-, TCG- and TT-compatible spatial coordinates: all frames
are treated as sharing one length unit, which is consistent with \lupnt{}'s
TDB-compatible convention but means that a state exchanged with a
TCB-compatible product needs an external $1-L_B$ scaling. There is no
station-displacement model (solid Earth tides, ocean loading, pole tide), so
an ITRF ground-station coordinate is a mean position. Nutation offsets
$\dd X, \dd Y$ from the EOP table are read but not applied to
$\mat{R}_{\mathrm{PN}}$.

\subsection{Expected accuracy}

\Cref{tab:fc-accuracy} summarises the differences measured against Orekit
v13.1 and GMAT R2026a on identical inputs and with \lupnt{}'s own physical
constants installed on both sides, so the residuals are algorithmic and
data-vintage differences rather than constant mismatches. Full methodology,
per-case numbers and the calibrated test tolerances are in
\file{docs/pages/cross_validation.rst} and the fixtures under
\file{cpp/test/orekit/} and \file{cpp/test/gmat/}.

\begin{table}[H]
  \centering
  \caption{Measured frame-conversion differences against external tools
    (June 2026 fixtures).}
  \label{tab:fc-accuracy}
  \small
  \begin{tabularx}{\textwidth}{>{\raggedright\arraybackslash}p{2.1cm} L L L}
    \toprule
    Conversion & vs Orekit (max) & vs GMAT (max) & Dominant cause \\
    \midrule
    GCRF $\leftrightarrow$ EME
      & \SI{1.8e-4}{\meter} (trans-lunar), $\sim\num{5e-13}\lVert\vec{r}\rVert$
      & \SI{54}{\meter} (trans-lunar), $\sim\num{1.6e-7}\lVert\vec{r}\rVert$
      & \lupnt{} and Orekit use the same IERS bias matrix; GMAT's
        ICRF$\leftrightarrow$MJ2000Eq differs by $\sim\num{1e-7}$~rad \\
    GCRF $\leftrightarrow$ ITRF
      & \SI{4.4}{\kilo\meter} trans-lunar ($\sim\num{1.3e-5}\lVert\vec{r}\rVert$); \SIrange{14}{98}{\meter} in LEO
      & \SI{1.6}{\kilo\meter} trans-lunar; \SIrange{34}{36}{\meter} in LEO
      & EOP vintage: a tens-of-ms UT1 difference is a
        $\num{1e-6}$--$\num{1e-5}$~rad polar rotation, scaling with radius \\
    GCRF $\to$ MOON\_CI
      & \SI{0.33}{\meter}
      & distribution level ($<\SI{200}{\meter}$)
      & pure DE440 translation; \lupnt{}, Orekit and GMAT R2026a all read DE440
        (GMAT through \lupnt{}'s \file{de440.bsp} via SPICE), so every column is
        a distribution/interpolation residual, not a DE-version offset \\
    Moon body-fixed
      & MOON\_ME vs IAU pole model: \SIrange{57}{200}{\meter} ($\sim\num{3e-5}\lVert\vec{r}\rVert$)
      & MOON\_PA vs SPICE Luna: distribution level ($<\SI{300}{\meter}$)
      & Orekit's analytic IAU pole approximates the DE mean-Earth axes to
        $\sim\num{3e-5}$~rad; GMAT's SPICE Luna PA kernel and \lupnt{}'s
        MOON\_PA, both on DE440, agree sub-arcsecond in orientation and at the
        distribution level in position \\
    Earth$\to$Moon, Earth$\to$Sun ephemeris
      & \SI{0.33}{\meter} / \SI{7.2}{\meter} (both $\sim\num{5e-11}$ relative)
      & ---
      & same DE440 data, independently distributed and interpolated \\
    \bottomrule
  \end{tabularx}
\end{table}

Two readings of \cref{tab:fc-accuracy} are worth stating explicitly. First,
the frames whose transformation is \emph{definitional} --- the frame bias,
the DE translation --- agree with an independent implementation at the
$\num{1e-13}$--$\num{1e-11}$ relative level, which is the floating-point floor.
Second, every metre-level and kilometre-level entry traces back to
Earth-orientation \emph{data}, not to the algorithm: the UT1 and
polar-motion table vintages. The DE ephemeris version is no longer a factor
now that \lupnt{}, Orekit and GMAT R2026a all read DE440, which is why the
Moon-centred columns have fallen to the ephemeris distribution level. That is
precisely the gap
\cpp{InitFrameConversionFromSpice} (\cref{sec:fc-spicefit}) exists to close.
For an application that needs SPICE-level Earth or lunar orientation over a
known window, fitting once at initialisation costs a few seconds and removes
the table-vintage term entirely; for everything else, the native analytic
path is self-contained, differentiable, and free of any runtime kernel
dependency.
